\documentclass[10pt,a4paper]{amsart}
\usepackage{amsmath,amssymb,amsthm,graphicx}
\usepackage{verbatim}
\newenvironment{reference}{\comment}{\endcomment}
%\newenvironment{reference}{}{}
\newenvironment{slogan}{\comment}{\endcomment}
\newenvironment{history}{\comment}{\endcomment}

% For commutative diagrams we use Xy-pic
\usepackage[all]{xy}

% We use 2cell for 2-commutative diagrams.
\xyoption{2cell}
% Theorem environments.
%
\theoremstyle{plain}
\newtheorem{theorem}[subsection]{Theorem}
\newtheorem{proposition}[subsection]{Proposition}
\newtheorem{lemma}[subsection]{Lemma}

\theoremstyle{definition}
\newtheorem{definition}[subsection]{Definition}
\newtheorem{example}[subsection]{Example}
\newtheorem{exercise}[subsection]{Exercise}
\newtheorem{situation}[subsection]{Situation}

\theoremstyle{remark}
\newtheorem{remark}[subsection]{Remark}
\newtheorem{remarks}[subsection]{Remarks}

\numberwithin{equation}{subsection}

% Macros
%
\def\lim{\mathop{\mathrm{lim}}\nolimits}
\def\colim{\mathop{\mathrm{colim}}\nolimits}
\def\Spec{\mathop{\mathrm{Spec}}}
\def\Hom{\mathop{\mathrm{Hom}}\nolimits}
\def\Ext{\mathop{\mathrm{Ext}}\nolimits}
\def\SheafHom{\mathop{\mathcal{H}\!\mathit{om}}\nolimits}
\def\SheafExt{\mathop{\mathcal{E}\!\mathit{xt}}\nolimits}
\def\Sch{\mathit{Sch}}
\def\Mor{\mathop{\mathrm{Mor}}\nolimits}
\def\Ob{\mathop{\mathrm{Ob}}\nolimits}
\def\Sh{\mathop{\mathit{Sh}}\nolimits}
\def\NL{\mathop{N\!L}\nolimits}
\def\CH{\mathop{\mathrm{CH}}\nolimits}
\def\proetale{{pro\text{-}\acute{e}tale}}
\def\etale{{\acute{e}tale}}
\def\QCoh{\mathit{QCoh}}
\def\Ker{\mathop{\mathrm{Ker}}}
\def\Im{\mathop{\mathrm{Im}}}
\def\Coker{\mathop{\mathrm{Coker}}}
\def\Coim{\mathop{\mathrm{Coim}}}

% Boxtimes
%
\DeclareMathSymbol{\boxtimes}{\mathbin}{AMSa}{"02}

%
% Macros for moduli stacks/spaces
%
\def\QCohstack{\mathcal{QC}\!\mathit{oh}}
\def\Cohstack{\mathcal{C}\!\mathit{oh}}
\def\Spacesstack{\mathcal{S}\!\mathit{paces}}
\def\Quotfunctor{\mathrm{Quot}}
\def\Hilbfunctor{\mathrm{Hilb}}
\def\Curvesstack{\mathcal{C}\!\mathit{urves}}
\def\Polarizedstack{\mathcal{P}\!\mathit{olarized}}
\def\Complexesstack{\mathcal{C}\!\mathit{omplexes}}
% \Pic is the operator that assigns to X its picard group, usage \Pic(X)
% \Picardstack_{X/B} denotes the Picard stack of X over B
% \Picardfunctor_{X/B} denotes the Picard functor of X over B
\def\Pic{\mathop{\mathrm{Pic}}\nolimits}
\def\Picardstack{\mathcal{P}\!\mathit{ic}}
\def\Picardfunctor{\mathrm{Pic}}
\def\Deformationcategory{\mathcal{D}\!\mathit{ef}}
\usepackage{cleveref}
\crefname{originaldefinition}{original definition}{original definitions}
\Crefname{originaldefinition}{Original definition}{Original definitions}
\makeatletter
\newlabel{homology-definition-filtered-complex}{{012L}{source}{}{sourcepointer-1}{}}
\newlabel{homology-definition-filtered-complex@cref}{{[originaldefinition][0][]012L}{[1][0][]source}{}{}{}}
\makeatother
\crefname{originaltheorem}{original theorem}{original theorems}
\Crefname{originaltheorem}{Original theorem}{Original theorems}
\makeatletter
\newlabel{theorem-K-injective-embedding-grothendieck}{{079P}{source}{}{sourcepointer-2}{}}
\newlabel{theorem-K-injective-embedding-grothendieck@cref}{{[originaltheorem][0][]079P}{[1][0][]source}{}{}{}}
\makeatother
\crefname{originallemma}{original lemma}{original lemmas}
\Crefname{originallemma}{Original lemma}{Original lemmas}
\makeatletter
\newlabel{derived-lemma-make-injective}{{013N}{source}{}{sourcepointer-3}{}}
\newlabel{derived-lemma-make-injective@cref}{{[originallemma][0][]013N}{[1][0][]source}{}{}{}}
\makeatother
\crefname{originallemma}{original lemma}{original lemmas}
\Crefname{originallemma}{Original lemma}{Original lemmas}
\makeatletter
\newlabel{lemma-derived-products}{{07D9}{source}{}{sourcepointer-4}{}}
\newlabel{lemma-derived-products@cref}{{[originallemma][0][]07D9}{[1][0][]source}{}{}{}}
\makeatother
\title{Realising an inverse system of derived objects by a filtered complex}
\author{The Stacks Project Authors}
\date{stacks-project:injectives.tex:local-snapshot-20261012; GNU FDL 1.2 or later}
\begin{document}
\maketitle
\noindent\textit{Editorial note.} All delivered mathematical text is the author's own, copied from the pinned Stacks Project chapter file sites-cohomology.tex; no statement or proof step was rewritten, reordered or removed.\par\medskip
\noindent\textit{Editorial note.} The proof is the author's complete construction; the single omitted step is the sentence 'We omit the verification that all the arrows in the diagram are termwise injective.', a routine check of the arrows of a displayed diagram whose definition is given in the lines immediately before it. It is preserved verbatim, and the rest of the argument, including the homotopy lifting, is written out in the counted span.\par\medskip
\noindent\textit{Editorial note.} Every labelled reference occurring inside the counted spans is carried in assembly.external\_references as a non-load-bearing navigation pointer (same-chapter ones at their definition site, cross-chapter ones at the invocation inside the counted span), so no load-bearing mathematics is replaced by a stub. No bibliography citation occurs inside the counted spans.\par\medskip
\noindent\textit{Editorial note.} The same-name check for this card scanned every staged Stacks chapter of this repository -- all sources/stacks-*/pinned-original-source.tex files, the card's own chapter excluded -- and the labels found are recorded in dedup.same\_name\_cross\_chapter on the card.\par\medskip
\noindent\textit{Editorial note.} The preamble is the chapter's own retained minimal preamble (sources/stacks-sites-cohomology/retained-preamble.tex), a verbatim excerpt of the native preamble: the Xy-pic setup and the native theorem declarations with the mathematical macros.\par\medskip
\noindent\textit{Allowed standard prerequisites (editorial).}\begin{itemize}
\item K-injective embeddings in a Grothendieck abelian category. Same chapter, Theorem theorem-K-injective-embedding-grothendieck and Theorem theorem-injective-embedding-grothendieck. Supply the K-injective complexes with injective terms that represent the given objects and the termwise injective maps of complexes into K-injective acyclic complexes used in the two constructions.
\item Products of K-injective complexes. Same chapter, Lemma lemma-derived-products. Supplies that products in the derived category are computed by products of K-injective complexes, which identifies the diagram built in the proof with the given one.
\item Filtered complexes. Homology, Definition homology-definition-filtered-complex. Supplies the notion of a filtered complex and of its filtration steps that the conclusion is about.
\item Termwise injective replacements. Derived Categories, Lemma derived-lemma-make-injective. Supplies the termwise injective maps of complexes representing the maps of the inverse system in the negative direction.
\end{itemize}\medskip
\ifdefined\hypertarget\hypertarget{sourcepointer-1}{}\fi
\noindent\textit{Source reference.} Original definition 012L: sources/stacks-homology/pinned-original-source.tex, source lines 6175--6175. Navigation pointer invoked at proof line 2611; the counted span contains the pointer only and the cited result is not reproduced inside it, so the pointer is registered as a non-load-bearing navigation stub in the established form for Stacks material. No bibliography citation is involved.\par\medskip
\ifdefined\hypertarget\hypertarget{sourcepointer-2}{}\fi
\noindent\textit{Source reference.} Original theorem 079P: sources/stacks-injectives/pinned-original-source.tex, source lines 2024--2024. Navigation pointer invoked at proof line 2552; the counted span contains the pointer only and the cited result is not reproduced inside it, so the pointer is registered as a non-load-bearing navigation stub in the established form for Stacks material. No bibliography citation is involved.\par\medskip
\ifdefined\hypertarget\hypertarget{sourcepointer-3}{}\fi
\noindent\textit{Source reference.} Original lemma 013N: sources/stacks-derived/pinned-original-source.tex, source lines 2840--2840. Navigation pointer invoked at proof line 2632; the counted span contains the pointer only and the cited result is not reproduced inside it, so the pointer is registered as a non-load-bearing navigation stub in the established form for Stacks material. No bibliography citation is involved.\par\medskip
\ifdefined\hypertarget\hypertarget{sourcepointer-4}{}\fi
\noindent\textit{Source reference.} Original lemma 07D9: sources/stacks-injectives/pinned-original-source.tex, source lines 2226--2226. Navigation pointer invoked at proof line 2699; the counted span contains the pointer only and the cited result is not reproduced inside it, so the pointer is registered as a non-load-bearing navigation stub in the established form for Stacks material. No bibliography citation is involved.\par\medskip
\begin{lemma}
\label{lemma-represent-by-filtered-complex}
Let $\mathcal{A}$ be a Grothendieck abelian category. Suppose given an object
$E \in D(\mathcal{A})$ and an inverse system $\{E^i\}_{i \in \mathbf{Z}}$
of objects of $D(\mathcal{A})$ over $\mathbf{Z}$ together with
a compatible system of maps $E^i \to E$. Picture:
$$
\ldots \to E^{i + 1} \to E^i \to E^{i - 1} \to \ldots \to E
$$
Then there exists a filtered complex $K^\bullet$ of $\mathcal{A}$
(Homology, Definition \ref{homology-definition-filtered-complex})
such that $K^\bullet$ represents $E$
and $F^iK^\bullet$ represents $E^i$ compatibly with the given maps.
\end{lemma}

\begin{proof}
By Theorem \ref{theorem-K-injective-embedding-grothendieck}
we can choose a K-injective complex $I^\bullet$
representing $E$ all of whose terms $I^n$ are injective
objects of $\mathcal{A}$.
Choose a complex $G^{0, \bullet}$ representing $E^0$.
Choose a map of complexes $\varphi^0 : G^{0, \bullet} \to I^\bullet$
representing $E^0 \to E$.
For $i > 0$ we inductively represent $E^i \to E^{i - 1}$
by a map of complexes
$\delta : G^{i, \bullet} \to G^{i - 1, \bullet}$
and we set $\varphi^i = \delta \circ \varphi^{i - 1}$.
For $i < 0$ we inductively represent $E^{i + 1} \to E^i$
by a termwise injective map of complexes
$\delta : G^{i + 1, \bullet} \to G^{i, \bullet}$
(for example you can use
Derived Categories, Lemma \ref{derived-lemma-make-injective}).
Claim: we can find a map of complexes
$\varphi^i : G^{i, \bullet} \to I^\bullet$
representing the map $E^i \to E$ and
fitting into the commutative diagram
$$
\xymatrix{
G^{i + 1, \bullet} \ar[r]_\delta \ar[d]_{\varphi^{i + 1}} &
G^{i, \bullet} \ar[ld]^{\varphi^i} \\
I^\bullet
}
$$
Namely, we first choose any map of complexes
$\varphi : G^{i, \bullet} \to I^\bullet$
representing the map
$E^i \to E$. Then we see that $\varphi \circ \delta$
and $\varphi^{i + 1}$ are homotopic by some homotopy
$h^p : G^{i + 1, p} \to I^{p - 1}$.
Since the terms of
$I^\bullet$ are injective and since $\delta$
is termwise injective, we can lift $h^p$ to
$(h')^p : G^{i, p} \to I^{p - 1}$.
Then we set $\varphi^i = \varphi + h' \circ d + d \circ h'$
and we get what we claimed.

\medskip\noindent
Next, we choose for every $i$ a termwise injective map of complexes
$a^i : G^{i, \bullet} \to J^{i, \bullet}$ with $J^{i, \bullet}$
acyclic, K-injective, with $J^{i, p}$ injective objects of $\mathcal{A}$.
To do this first map $G^{i, \bullet}$ to the cone on the identity
and then apply the theorem cited above.
Arguing as above we can find maps of complexes
$\delta' : J^{i, \bullet} \to J^{i - 1, \bullet}$ such that the diagrams
$$
\xymatrix{
G^{i, \bullet} \ar[r]_\delta \ar[d]_{a^i} &
G^{i - 1, \bullet} \ar[d]^{a^{i - 1}} \\
J^{i, \bullet} \ar[r]^{\delta'} &
J^{i - 1, \bullet}
}
$$
commute. (You could also use the functoriality of cones plus the
functoriality in the theorem to get this.)
Then we consider the maps
$$
\xymatrix{
G^{i + 1, \bullet} \times \prod\nolimits_{p > i + 1} J^{p, \bullet}
\ar[r] \ar[rd] &
G^{i, \bullet} \times \prod\nolimits_{p > i} J^{p, \bullet}
\ar[r] \ar[d] &
G^{i - 1, \bullet} \times \prod\nolimits_{p > i - 1} J^{p, \bullet}
\ar[ld] \\
& I^\bullet \times \prod\nolimits_p J^{p, \bullet}
}
$$
Here the arrows on $J^{p, \bullet}$ are the obvious ones
(identity or zero). On the factor $G^{i, \bullet}$ we use
$\delta : G^{i, \bullet} \to G^{i - 1, \bullet}$, the map
$\varphi^i : G^{i, \bullet} \to I^\bullet$, the zero map
$0 : G^{i, \bullet} \to J^{p, \bullet}$ for $p > i$, the map
$a^i : G^{i, \bullet} \to J^{p, \bullet}$ for $p = i$, and
$(\delta')^{i - p} \circ a^i = a^p \circ \delta^{i - p} :
G^{i, \bullet} \to J^{p, \bullet}$ for $p < i$.
We omit the verification that all the arrows
in the diagram are termwise injective. Thus we obtain a filtered
complex. Because products in $D(\mathcal{A})$ are given by
taking products of K-injective complexes
(Lemma \ref{lemma-derived-products})
and because $J^{p, \bullet}$ is zero in $D(\mathcal{A})$
we conclude this diagram represents the given diagram in the derived
category. This finishes the proof.
\end{proof}

\end{document}
